\honest{Church vs. Taskforce}{Eliezer Yudkowsky}{2009}

I have warned before against envying crazy groups and against writing ``hymns to the nonexistence of God''; a good atheistic hymn is simply a song about anything worth singing about that is not religious. But religion does fill certain holes in people's minds, some worth filling. Delete religion, and the largest gap would be the church, the community. How many people stay religious without quite believing, just to keep their neighbours? ``I would guess\ldots\ probably quite a lot.'' This is ``probably something I don't understand all that well, myself.'' Brownies and babysitting come to mind first. But the right question is what a hunter-gatherer band gives its people, and what modern life lacks; if a church fills only some of that, let us do better, without assuming we must meet on Sunday mornings under stained glass. \nb{Research published the year after the post supports the general point. Lim and Putnam (2010) found that religious people are more satisfied with their lives because they attend services and have friends in their congregation, with little evidence that private belief matters apart from that.}

To break the mold, note that some offices supply community too: brownies, babysitters, maybe help in a catastrophe. Not everyone is that lucky. But neither a church nor a workplace was designed to be a community. So you could suspect, though it ``would be better tested experimentally,'' that early Sunday mornings, formal clothes and weekly sermons on one theme are not optimal, that supporting a church and a pastor is expensive when many communities could share one building, that churches do too little matchmaking because they enforce medieval moralities, and that all of it should be tested. By ``optimal'' I mean optimal for a community built as a community; stained glass makes sense if you want religion. Walking past the churches of my city, I think ``These buildings look really, really expensive, and there are too many of them.'' Starting over, you might have one large building for the occasional wedding, shared by different communities at different times, with a big video screen for speakers, lectures or films. Stained glass would not be a high priority. \nb{Many church buildings are already shared. Partners for Sacred Places, an organization that supports historic religious buildings, reports that most people served by programs housed in congregations are not members, and that in one study of 90 congregations only 9 per cent of visits were for worship.}

Could the help a congregation gives in trouble be improved by an explicit rainy-day fund or insurance? Possibly not; bringing in explicit finance changes things oddly. Perhaps keeping insured should be a condition of membership. But churches provide community without admitting that this is ``nearly all of what people get out of it,'' and friendly workplaces provide it by accident. Think explicitly about giving people a band, and good ideas appear. Welcoming newcomers is not a sermon topic; right after a move is when someone most needs community, and it is a chance for the band to grow. Tribes might even compete at quarterly exhibitions to capture newcomers. \nb{Asked by Pew in 2018, eight in ten regular churchgoers called becoming closer to God a very important reason for going, and about one in twenty named community as the main reason. Stated reasons are not the same as benefits, but ``nearly all'' goes beyond both kinds of evidence.}

Can a community have no purpose beyond itself? Hunter-gatherer bands had one, survival. Anything people have in common, especially a goal, tends to define a community, so why not use that? But on the Internet supporters are too scattered; the only member of the Church of the Subgenius in your city gains little, and the Internet is not yet a substitute for physical presence. \nb{Lim and Putnam found that the benefit of friends in a congregation depended on a strong religious identity. Whether a shared secular goal can play that part is untested.}

So my proposal: rationalist communities as taskforces working on the world's problems, each built around the most specific interest that can support ``a decent-sized band.'' If your city lacks 50 fellow Linux programmers, settle for 15 open-source programmers, or at first 15 rationalists improving the world in their various ways. Such tasks need communities anyway, and the energy once spent on religious institutions could go there, with purposes admirable without delusion filling the gap. Colleges seem to support communities well, so learning is another possible center; with a large screen in the building, I would challenge the idea that adult learning must happen on distant, expensive campuses with teachers who would rather be doing something else. Is this a dream? ``Maybe. Probably.'' It could be built step by step in a large enough city with enough rationalists who have heard of it. As the churches fade, ``we don't need artificial churches, but we do need new idioms of community.''
